% =====================================================================
% P1 Appendix A.6 --- Theorem 4 (within-class non-identification without I5/I5$\prime$).
%
% Source of truth: drafts/proofs-props8-12.md v1.4, Section 6 (FS-11).
% This file is a fragment: it is meant to be \input under
% \appendix \section{Proofs} after A.5.
%
% Theorem-like environments use the explicit paper numbering
% (Theorem 4).  Equation tags (6.1)--(6.4) follow the frozen source;
% do not renumber without updating the main text.
% =====================================================================

\newtheorem*{aVIthm}{Theorem 4}

% Fixed equation tags: keep the frozen printed numbers while giving
% hyperref a unique destination for every display.
\providecommand{\theHequation}{}
\let\aVItheHequation\theHequation
\newcounter{aVIanchor}
\renewcommand{\theHequation}{aVI.\arabic{aVIanchor}}
\newcommand{\aVIfixedtag}[1]{\stepcounter{aVIanchor}\tag{#1}}

\subsection{Within-class non-identification without cross-event variation (Theorem~4)}\label{app:a6}

This appendix gives the finite-sample two-point lower bound behind
Theorem~4.  The population statement is the existence result summarized
in Section~5.3; the version below strengthens it to an arbitrary sample
size and an arbitrary observable-data estimator.

\paragraph{Model class.}
Consider the class $\mathcal M$ of models in which the freeze-period
interior path is unobserved, only OPEN-period data and noisy endpoints
of each freeze event are observed, and the open-period transition $T$,
the observation kernel $h$, the endpoint noise, and the calendar are
common.  The fixed behavioral and logging policy depends only on
observable data, but the freeze-period bridge law may differ across
models.  Suppose the bridge expectation of the reward $\rho$ is not
determined by the endpoints: there are two bridge laws $Q_1,Q_2$,
compatible with the same open transition, observation kernel, and
endpoint law, such that
\[
\mathbb E_{Q_1}[\rho]\ne\mathbb E_{Q_2}[\rho].
\]
The construction in \eqref{eq:a6-6.2} below, with the reward
$-\lambda\,\mathrm{QV}_{\mathrm{freeze}}$, satisfies this condition.

\begin{aVIthm}
There exist models $P,Q\in\mathcal M$, independent of $n$ and of the
estimator, such that for every $n\ge1$ and every estimator $\hat V_n$ that
depends only on observable data (and possibly on external randomization):
\begin{enumerate}
\item[(i)] the joint law of the observable data is identical under $P$
and $Q$, for every $n$;
\item[(ii)] $|V^\pi(P)-V^\pi(Q)|=\Delta>0$;
\item[(iii)]
\begin{equation}\label{eq:a6-6.1}
\max\Bigl(
\mathbb E_P\bigl|\hat V_n-V^\pi(P)\bigr|,\,
\mathbb E_Q\bigl|\hat V_n-V^\pi(Q)\bigr|
\Bigr)
\ge\frac{\Delta}{2}.
\aVIfixedtag{6.1}
\end{equation}
\end{enumerate}
Consequently, no estimator is uniformly consistent over
$\mathcal M$; the worst-case risk of any estimator is at least half the
value gap.
\end{aVIthm}

\paragraph{Construction: Brownian bridge and quadratic-variation reward.}
Fix one freeze event of duration $\tau>0$ and endpoints $s_-,s_+$.
Take two bridge laws
\begin{equation}\label{eq:a6-6.2}
\begin{aligned}
Q_\sigma:\quad & \log S \text{ is a Brownian bridge with volatility }
\sigma \text{ given the endpoints},\\
Q_{\sigma'}:\quad & \text{same endpoints and same open transition,
volatility } \sigma'\ne\sigma.
\end{aligned}
\aVIfixedtag{6.2}
\end{equation}
The endpoint marginals are the same because the endpoints are fixed;
the open transition $T$, observation kernel $h$, and endpoint noise are
the same; and the freeze-period interior is unobserved.  Hence every
observable has the same joint law under the two models, for every $n$.
Let $P$ be the full law with $Q_\sigma$ and $Q$ the full law with
$Q_{\sigma'}$, matching the notation of Theorem~4.

Let the reward contain the term $-\lambda\,\mathrm{QV}_{\mathrm{freeze}}$
with $\lambda>0$, where
\[
\mathrm{QV}_{\mathrm{freeze}}
=\int_{\mathrm{freeze}}(d\log S_t)^2
\]
is the quadratic variation of the log price over the freeze period.  For
a Brownian bridge with fixed endpoints,
\begin{equation}\label{eq:a6-6.3}
\mathbb E_{Q_\sigma}[\mathrm{QV}_{\mathrm{freeze}}]=\sigma^2\tau,
\qquad
\mathbb E_{Q_{\sigma'}}[\mathrm{QV}_{\mathrm{freeze}}]
=\sigma'^2\tau.
\aVIfixedtag{6.3}
\end{equation}
The linear interpolation term has zero quadratic variation and zero
cross-variation with the Brownian bridge, and the standard bridge has
quadratic-variation expectation $\tau$.  Therefore
\begin{equation}\label{eq:a6-6.4}
V^\pi(P)-V^\pi(Q)
=-\lambda\tau(\sigma^2-\sigma'^2)
=\lambda\tau(\sigma'^2-\sigma^2).
\aVIfixedtag{6.4}
\end{equation}
Taking $\sigma'\ne\sigma$ gives
$\Delta=\lambda\tau|\sigma'^2-\sigma^2|>0$.

If the reward depends only on the endpoints and has no path term, the two
bridge laws give the same $V^\pi$.  The path-dependence condition is
therefore essential, not a technical detail.  Quadratic variation is
only a convenient sufficient example: any path functional whose
expectation differs under $Q_\sigma$ and $Q_{\sigma'}$ can replace it.

\begin{proof}
\begin{enumerate}
\item \emph{The observable law is the same.}  By construction, $P$ and
$Q$ differ only in the freeze-period interior bridge law.  All
observation channels (OPEN data, endpoints, calendar, and actions) do
not touch the interior path, and the behavioral policy depends only on
observables.  Hence the observable data $O_n$ has the same joint law
under $P$ and $Q$.

\item \emph{The estimator has the same distribution.}  The estimator
$\hat V_n(O_n)$ depends only on $O_n$, whose law is the same under $P$
and $Q$; hence $\hat V_n$ has the same distribution.  If the estimator
is randomized, couple the same random seed under the two laws.

\item \emph{Two-point bound.}  Let $Z=\hat V_n$ under the common law,
$V_P=V^\pi(P)$, and $V_Q=V^\pi(Q)$.  The pointwise triangle inequality
gives
\[
|Z-V_P|+|Z-V_Q|\ge|V_P-V_Q|=\Delta.
\]
Taking expectations under the common law (if the expectations are
infinite, the conclusion is trivial) yields
\[
\mathbb E_P|\hat V_n-V_P|+\mathbb E_Q|\hat V_n-V_Q|\ge\Delta,
\]
so the maximum of the two terms is at least $\Delta/2$.  For squared
loss, the inequality
\[
\Delta^2\le2\bigl(|Z-V_P|^2+|Z-V_Q|^2\bigr)
\]
gives
\[
\max\Bigl(
\mathbb E_P(\hat V_n-V_P)^2,\,
\mathbb E_Q(\hat V_n-V_Q)^2
\Bigr)
\ge\frac{\Delta^2}{4}.
\]
\end{enumerate}
\end{proof}

\paragraph{Relation to the population statement.}
The population proposition asserts that two observationally equivalent
models can have different values, so the value is not identified.
Theorem~4 strengthens this in three ways: it holds for every finite $n$,
it gives the explicit two-point risk bound $\Delta/2$, and it isolates
the path-dependence condition with the quadratic-variation construction.
It does not prove that a particular estimator, such as a plug-in
estimator, fails on the observed data; it proves that no estimator is
uniformly consistent over the model class.  This is an
information-theoretic statement, distinct from the empirical
correspondence in the stress variant, which tests the unfreeze-jump
mechanism rather than the lower bound.

\paragraph{Scope.}
\begin{itemize}
\item The theorem does not claim that every bridge law leads to a
different $V^\pi$; an endpoint-only reward is a counterexample.
\item It does not claim that $\Delta$ is large in real data; only that
$\Delta>0$.
\item It does not claim that quadratic variation is the only mechanism.
Occupation time, residence time, and other path functionals with
different expectations under the two bridge laws work equally well.
\item It does not claim that every volatility is unidentified in the
population model.  The result is an existence statement: after removing
I5/I5$'$, the value is not identified.
\end{itemize}

\let\theHequation\aVItheHequation
